\section{Evaluación y síntesis de resultados}

Para cuantificar el impacto real de la migración hacia arquitecturas EDA, no nos basamos en promedios abstractos. Diseñamos un esquema de evaluación que clasifica los resultados reportados en los 20 estudios bajo tres indicadores de rendimiento clave (KPIs): \textbf{Precisión de Detección}, \textbf{Latencia de Procesamiento} y \textbf{Escalabilidad Horizontal}.

A diferencia de las pruebas de laboratorio controladas, estos datos provienen de implementaciones en entornos reales, lo que les otorga un peso estadístico significativo sobre la viabilidad de estas tecnologías en producción.

\subsection{Impacto de la Inteligencia Contextual}
El dato más revelador no proviene de la velocidad pura, sino de la capacidad del sistema para adaptarse al contexto local. Perera et al. \cite{perera2024traffic} demostraron que utilizar algoritmos de reconocimiento genéricos es insuficiente para escenarios complejos. Aunque su modelo base (YOLOv5m) alcanzó una precisión media (mAP) del 96.1\% en condiciones ideales, la lectura de matrículas en el mundo real caía drásticamente sin lógica de negocio adicional.

Como se observa en la \cref{tab:mejora_precision}, la introducción de un algoritmo \textit{ad-hoc} que "entiende" el formato de las matrículas de Sri Lanka aumentó la tasa de éxito en 30 puntos porcentuales. Esto valida que la capa semántica es tan crítica como el modelo de IA: el algoritmo no solo debe "ver" caracteres, debe saber qué esperar.

\begin{table}[htbp]
\centering
\caption{Comparación de precisión en detección de matrículas (Sri Lanka)}
\label{tab:mejora_precision}
\resizebox{\columnwidth}{!}{%
\begin{tabular}{lcc}
\toprule
\textbf{Método de Procesamiento} & \textbf{Detecciones} & \textbf{Precisión (\%)} \\
\midrule
Solo Preproc. + OCR & 27 / 50 & 54\% \\
Preproc. + Algoritmo Ad-hoc + OCR & 42 / 50 & \textbf{84\%} \\
\bottomrule
\end{tabular}%
}
\end{table}

Este hallazgo se alinea con el estudio de Coretti, Varile y Bertaina \cite{coretti2025neuromorphic}, donde el cambio a visión neuromórfica (procesar eventos de luz en lugar de cuadros completos) permitió mantener la precisión operativa reduciendo el ruido de fondo, algo vital para la detección de objetos tenues como la basura espacial.

\subsection{Eficacia de los Modelos de Visión en el Borde}
Profundizando en la detección, los resultados de Kul y Tashiev \cite{kul2021microservices} desmienten el mito de que "más grande es mejor". Al comparar variantes de redes neuronales, descubrieron que el modelo ligero \textbf{YOLOv4-tiny} alcanzó una Precisión Promedio (AP) del 98.39\% en la clasificación de vehículos, superando paradójicamente a modelos más pesados en el dataset de prueba. 

Además, su análisis sobre la detección de color arrojó una lección técnica importante para el preprocesamiento en el borde: el uso del espacio de color HSV (Hue, Saturation, Value) logró una precisión del 95.63\%, superior al estándar RGB (95.29\%). Esto confirma que transformaciones simples de datos en el origen pueden mejorar la inferencia sin costo computacional adicional.

\subsection{Desbloqueo del Rendimiento (Throughput)}
En términos de velocidad, la evidencia apunta a que el cuello de botella tradicional no era el procesamiento, sino la capacidad de ingesta de datos. El sistema legado de Tecgraf/PUC-Rio \cite{laigner2020monolithic} sufría bloqueos constantes por concurrencia al intentar escribir en bases de datos compartidas.

Para probar la escalabilidad de la solución EDA, Kul y Tashiev \cite{kul2021microservices} sometieron su arquitectura basada en Apache Kafka a pruebas de carga incremental. Los resultados, detallados en la \cref{tab:kafka_performance}, son contundentes: el sistema no solo soportó un aumento exponencial en la carga (de 10 a 100,000 mensajes), sino que el tiempo de consumo se mantuvo extremadamente eficiente. Notablemente, procesar 100,000 eventos tomó apenas 1.2 segundos (1209 ms), demostrando que la arquitectura asíncrona puede manejar picos de tráfico masivos que colapsarían a un servidor HTTP tradicional.

\begin{table}[htbp]
\centering
\caption{Rendimiento de Apache Kafka (Producción vs. Consumo)}
\label{tab:kafka_performance}
\resizebox{\columnwidth}{!}{%
\begin{tabular}{rcc}
\toprule
\textbf{Carga (Mensajes)} & \textbf{T. Producción (ms)} & \textbf{T. Consumo (ms)} \\
\midrule
10 & 40 & 54 \\
100 & 145 & 259 \\
1,000 & 285 & 295 \\
10,000 & 400 & 456 \\
100,000 & 2,863 & \textbf{1,209} \\
\bottomrule
\end{tabular}%
}
\end{table}

\subsection{Eficiencia Operativa y Reducción de Carga Humana}
Finalmente, los resultados cualitativos son igual de impactantes. Bruns y Dunkel \cite{bruns2014fusion} reportaron que, al implementar Procesamiento de Eventos Complejos (CEP) en la coordinación de ambulancias, se eliminó la necesidad de que los conductores reportaran manualmente sus estados ("en ruta", "libre"). El sistema infiere el estado automáticamente cruzando eventos de GPS y telemetría, reduciendo el error humano y acelerando los tiempos de despacho.

De manera similar, la reorganización lógica de los sensores mediante árboles de acceso (SAT) propuesta por Shujaat et al. \cite{shujaat2021improvements} mejoró la eficiencia de recuperación de datos en entornos IoT masivos, demostrando que la optimización no siempre requiere hardware nuevo, sino estructuras de datos más inteligentes.

En resumen, los datos confirman nuestra hipótesis inicial: la mejora es estructural. Los sistemas que adoptaron EDA y filtrado en el borde \cite{salamon2018gateway} lograron procesar volúmenes de datos superiores con mayor precisión, validando que la asincronía es el único camino viable para la escalabilidad en el IoT moderno.